% Pista ana-26s-q4a · Anàlisi
% Procedència: PAU Catalunya 2026 · Convocatòria de setembre · Sèrie 2 · Exercici 4 — Opció A
\documentclass[11pt,a4paper]{article}
\usepackage{pista}
\pistainfo{Catalunya 2026}{Convocatòria de setembre}{Sèrie 2}

\begin{document}

\section*{\textcolor{blauPista}{Exercici 4 --- Opció A}}

\noindent L'Aina vol construir una capsa de cartró amb forma de prisma rectangular amb tapa, per guardar-hi la tauleta i el carregador.

\noindent L'Aina construirà la capsa a partir d'un cartró rectangular de $80\,\text{cm} \times 50\,\text{cm}$. A la figura següent es mostra el desplegament de la capsa; les zones ratllades són les parts que l'Aina retallarà i retirarà per tal de poder-la muntar:

\noindent\textit{(A la figura: el costat de $80\,\text{cm}$ és l'horitzontal. D'esquerra a dreta hi ha una franja d'amplada $x$, un rectangle, una altra franja d'amplada $x$ i un altre rectangle. Les zones ratllades són els quadrats de costat $x$ de dalt i de baix de les dues franges i, al darrer rectangle, les tires de dalt i de baix d'amplada $x$.)}

\subsection*{4.1. \normalfont Comproveu que el volum de la capsa resultant és $V(x) = 2x^{3} - 130x^{2} + 2\,000x$, on $x$ és la longitud del costat del quadrat que cal retallar dels cantons per a poder-la muntar. \hspace*{\fill}\textnormal{[1~punt]}}

\pista{Prova-ho en primer lloc amb un número inventat. Per exemple, si $x = 5$, al llarg dels $80\,\text{cm}$ hi ha dues franges de $5\,\text{cm}$, i en queden $80 - 10 = 70\,\text{cm}$ per a la base i la tapa. Com que la base i la tapa són iguals, la base fa $\frac{70}{2} = 35\,\text{cm}$ de llargada. En vertical, a dalt i a baix de la base hi ha els altres dos laterals, i per tant l'amplada de la base fa $50 - 2 \cdot 5 = 40\,\text{cm}$. L'altura de la capsa és $5\,\text{cm}$.

\medskip
Ara fes el mateix amb $x$ en lloc de $5$. Finalment, multiplica la llargada, l'amplada i l'altura, i desenvolupa el producte fins a arribar a l'expressió de l'enunciat.}

\subsection*{4.2. \normalfont Calculeu les dimensions de la capsa per tal que el seu volum sigui màxim. Quin és aquest volum màxim? \hspace*{\fill}\textnormal{[1,5~punts]}}

\pista{En primer lloc, deriva $V(x)$ i resol $V'(x) = 0$. Et quedarà una equació de segon grau, més senzilla si abans la divideixes per $2$. Obtindràs dues solucions: queda't amb la que té sentit, és a dir, aquella amb la qual les tres dimensions de l'apartat 4.1 són positives.

\medskip
A continuació, comprova que és un màxim fent servir la segona derivada $V''$.

\medskip
Finalment, tingues en compte que la teva resposta ha d'incloure, d'una banda, les tres dimensions de la capsa; i d'altra banda, el volum màxim, amb les unitats.}

\end{document}
